\section{Método}

\subsection{Arquitecturas}
\label{sec:arquitecturas}

Ambos modelos comparten el bloque de GPT-2 \citep{radford2019language}:
normalización previa, atención multicabeza causal, red \emph{feed-forward} de
dimensión $4d$ con GELU y conexiones residuales,
$B_\theta(\hid{}) = \mathbf{u} + \FFN(\LN_2(\mathbf{u}))$ con
$\mathbf{u} = \hid{} + \MHA(\LN_1(\hid{}))$. La entrada es la suma de los
\emph{embeddings} de carácter y de posición (aprendido), $\hid{0}$, y la salida
$\mathrm{softmax}(W_o\,\LN(\hid{D}))$ con $D$ la profundidad efectiva. Un
\textbf{GPT} de $L$ capas usa bloques distintos, $\hid{\ell}=B_{\theta_\ell}(\hid{\ell-1})$;
un \textbf{Looped Transformer} de $N$ iteraciones reutiliza el mismo,
\begin{equation}
\hid{t} = B_\theta(\hid{t-1}),\qquad t = 1,\dots,N,
\end{equation}
con $D=L$ y $D=N$ respectivamente. Con dimensión $d$, vocabulario $V$ y contexto
$T$, un bloque tiene $P_B = 12d^2 + 10d$ parámetros y el resto del modelo
$P_E = (2V+T)d + 2d$, de modo que $P_\text{GPT}(L) = L\,P_B + P_E$ y
$P_\text{Looped}(N) = P_B + P_E$ no depende de $N$, mientras que el cómputo por
token, $F(D) \approx D\,(24d^2 + 4Td) + 2dV$, crece con $D$ en ambos.

\subsection{Implementación y validación}
\label{sec:validacion}

Partimos del código de \citet{raschka2024build}, con la atención explícita para
devolver sus pesos; ambos modelos exponen la misma interfaz (logits, estados
ocultos y atenciones por paso). Comprobamos la corrección con pruebas
automáticas: Looped con $N{=}1$ coincide con GPT con $L{=}1$ y mismos
pesos ($10^{-6}$); Looped con $N\in\{2,4,8,12\}$ coincide, en logits y en todos
los estados ocultos, con un GPT de $N$ capas con copias del bloque compartido
($10^{-5}$); el número de parámetros no depende de $N$ y la atención es causal.

\subsection{Datos}
\label{sec:datos}

Cada ejemplo es una historia breve en castellano cuyo nivel es el número de
operaciones, de 1 a 6: suma, resta, doble, triple, mitad y tercera y cuarta
parte, con números de hasta 20 en el enunciado y respuestas entre 1 y 99. Hay
relatos con un protagonista o un recipiente (tipo~A) y problemas de comparación
entre varias personas (tipo~B) cuyas frases pueden estar desordenadas, en
estilo detallado o compacto y a partir de plantillas. Un ejemplo de nivel~2 con frases desordenadas:
\begin{quote}\small
\emph{Julia tiene la mitad que Laura. Laura tiene el triple de postales que Pablo.
Pablo tiene 18 postales. Entonces, ¿cuántas postales tiene Julia?}\ \ $\to$ 27
\end{quote}
Cada ejemplo guarda su traza de operaciones, con la que se verifica la
respuesta, y ninguna pregunta de validación o de test aparece en el
entrenamiento. Los niveles 1--4 se reparten en entrenamiento (200\,000
ejemplos, 50\,000 por nivel), validación (4\,000) y test en distribución
(ID, 8\,000, 2\,000 por nivel). Un segundo test de plantillas nuevas (8\,000
ejemplos) usa 18 redacciones reservadas que no aparecen en el entrenamiento.
Los niveles 5 y 6 (4\,000 ejemplos) solo se usan como test fuera de
distribución (OOD). Las heurísticas sin razonamiento (respuesta más frecuente, primer o último número) aciertan entre el 0.6 y el 7.6\,\%.

\subsection{Entrenamiento y evaluación}
\label{sec:entrenamiento}

El texto de entrada es \texttt{Pregunta: \dots\textbackslash nRespuesta: } seguido
del número y de un fin de secuencia, tokenizado por caracteres (97 símbolos); la
pérdida se calcula solo sobre la respuesta. Usamos AdamW
\citep{loshchilov2019decoupled} ($\beta=(0.9, 0.95)$, decaimiento 0.1), tasa
$10^{-3}$ con calentamiento de 200 pasos y decaimiento coseno hasta el 10\,\%,
lotes de 64, recorte de gradiente 1.0 y precisión mixta, durante 42\,000 pasos y con tres semillas (0, 1 y 2).
Evaluamos por generación voraz (hasta 6 caracteres, primer entero extraído) y
medimos la exactitud por nivel (coincidencia exacta); en los modelos con bucles variamos las
iteraciones en test, $N'$, de 1 a $2N$ sin reentrenar.
